% ============================================================
% PARTE FELIPE — presentación del paper y supuestos a fondo
% (el CAPM se explica acá, no en la parte de tasa de descuento)
% Nota: los títulos nunca llevan "Sección N"; las referencias a
% secciones del paper van solo en comentarios.
% ============================================================
% ---------- Sumario ----------
% (mapea el arco completo: 1-2 Felipe, 3 Elias, 4 Mariano, 5 Agustín)
\begin{frame}{Sumario}
  \vspace{10mm}
  \begin{enumerate}
    \item El error del VAN tradicional y el valor de esperar
    \item El problema de inversión y los supuestos del modelo
    \item La regla óptima y el valor de la opción de esperar
    \item La tasa de descuento (CAPM) y las extensiones del modelo
    \item Resultados numéricos, el monopolista y conclusiones
  \end{enumerate}
\end{frame}

\begin{frame}{Motivación: el error del VAN tradicional}
  \begin{itemize}
    \item Construir es \concepto{irreversible} --- el capital queda hundido,
          sin uso alternativo ---, mientras que diferir la inversión es
          \emph{reversible}. Esa asimetría es el corazón del problema.
    \item Invertir hoy no compite solo con ``no invertir'': compite con
          \emph{todas} las alternativas futuras, mutuamente excluyentes, de
          invertir más adelante.
  \end{itemize}
  \begin{cajatexto}
    Invertir apenas ``VAN $>0$'' (en $\rj{V} = \az{F}$) solo es óptimo en dos
    casos degenerados: si \az{$\sigma_v^2 = \sigma_f^2 = 0$} o si
    \rj{$\alpha_v \to -\infty$}. En cualquier otro caso hace falta un
    \textbf{margen positivo}.
  \end{cajatexto}
\end{frame}

\begin{frame}{Objetivo del paper}
  \vspace{5mm}
  El objetivo es estudiar el \concepto{timing óptimo} de la inversión en un
  proyecto irreversible, cuando los beneficios esperados y los costos de
  inversión siguen procesos estocásticos en tiempo continuo.
  \par\vspace{6mm}
  A diferencia de los modelos previos, deriva una \emph{regla de
  inversión óptima} y una fórmula explícita para valuar la opción de esperar,
  incorporando formalmente la aversión al riesgo de los
  inversores.
  \par\vspace{6mm}
  Y demuestra, mediante simulaciones numéricas, que el valor de esa opción es
  significativo: aplicar la regla tradicional del ``VAN $=0$'' genera
  \resultado{pérdidas sustanciales}.
\end{frame}

% (Sección II del paper) — V y F (MBG) + supuesto de vida infinita
\begin{frame}{El problema de inversión}
  \begin{itemize}
    \item En cualquier momento $t$, la firma puede pagar un costo \az{$F_t$}
          para instalar un proyecto cuyo valor presente de flujos netos
          esperados es \rj{$V_t$}.
    \item Ambos son estocásticos y siguen \concepto{movimientos brownianos
          geométricos}, con $dz_v$, $dz_f$ procesos de Wiener:
  \end{itemize}
  \begin{ec}
    \rj{\frac{dV}{V} = \alpha_v\,dt + \sigma_v\,dz_v}
    \qquad\qquad
    \az{\frac{dF}{F} = \alpha_f\,dt + \sigma_f\,dz_f}
  \end{ec}
  \vspace{-3mm}
  \begin{itemize}
    \item El MBG hace del cero una \emph{barrera absorbente} natural.
    \item \concepto{Vida infinita} (\az{$T=\infty$}): umbral óptimo constante.
    \item La inversión se hace \emph{toda de una vez}, no por etapas.
  \end{itemize}
\end{frame}

% Caso especial: costo fijo (a continuación de V y F)
\begin{frame}{Caso especial: costo de inversión fijo}
  \begin{itemize}
    \item Si el costo \az{$F$} es constante, su volatilidad y su deriva se
          anulan:
  \end{itemize}
  \begin{ecmed}
    \az{\sigma_f = 0}, \qquad \az{\alpha_f = 0}
  \end{ecmed}
  \begin{itemize}
    \item Lejos de perder la solución analítica, el problema de dos variables
          estocásticas colapsa a una sola (\rj{$V$}).
    \item Es el caso clásico de la \concepto{opción americana perpetua},
          resuelto con fórmula analítica cerrada.
  \end{itemize}
\end{frame}

% Vencimiento — incorpora el supuesto de salto de Poisson
% OJO: se solapa con las slides de Mariano (Saltos en Vt / obsolescencia) → NOTAS-CONFLICTOS.md
\begin{frame}{¿Y si la oportunidad puede vencer?}
  \begin{itemize}
    \item Con fecha de vencimiento fija ($T<\infty$) \resultado{no hay
          solución analítica}: cerca del límite, la frontera óptima de
          ejercicio cambia con el calendario y se necesitan métodos numéricos
          (árboles binomiales, diferencias finitas).
    \item Con \emph{vencimiento incierto} sí: se modela con un
          \concepto{salto de Poisson}, probabilidad constante \az{$\lambda\,dt$}
          de que \rj{$V$} caiga a cero (obsolescencia, entrada de un rival).
    \item La estructura analítica queda idéntica al caso base; el riesgo de
          salto solo penaliza la tasa de descuento:
  \end{itemize}
  \cajaec{\az{\mu} \;\longrightarrow\; \rj{\mu + \lambda}}
\end{frame}

\begin{frame}{Un mismo marco, varios problemas}
  \begin{itemize}
    \item \textbf{Monopolista}: tras instalar la capacidad queda protegido
          de la competencia; los flujos del proyecto están resguardados.
    \item \textbf{Industria competitiva}: rentas temporales --- la entrada
          rezagada de competidores empuja los beneficios hacia los costos
          (deriva post-inversión menor o negativa).
    \item \textbf{Intercambio de activos}: invertir equivale a cambiar un
          activo riesgoso (\az{$F$}) por otro (\rj{$V$}).
    \item \textbf{Desguace óptimo}: reinterpretando \az{$F$} como el valor
          en operación y \rj{$V$} como el valor de descarte.
  \end{itemize}
\end{frame}

% CAPM — supuesto de inversores diversificados + rol del CAPM en la tasa de descuento
% (viene del cuadro de supuestos, ahora desarmado)
% OJO: se solapa con la slide de CAPM de Mariano (derivación de μ) → NOTAS-CONFLICTOS.md
\begin{frame}{CAPM}
  \begin{itemize}
    \item Opciones y proyectos pertenecen a \concepto{portafolios bien
          diversificados}: la aversión al riesgo entra solo a través del
          \concepto{riesgo sistemático}.
    \item Ese supuesto habilita usar el CAPM para fijar la tasa de
          descuento correcta de los flujos: la tasa de rendimiento
          de equilibrio exigida a la inversión.
  \end{itemize}
  \begin{cajatexto}
    El CAPM es la pieza que permite obtener \az{$\mu$}, la tasa a la que se
    descuentan los flujos. Su desarrollo se retoma más adelante.
  \end{cajatexto}
\end{frame}

% Dividendo — δ = α̂v − αv ⟹ αv < α̂v
% (μ → α̂v por decisión del grupo: μ queda reservado para la tasa de la opción)
\begin{frame}{¿Por qué $\alpha_v < \hat{\alpha}_v$? La analogía del dividendo}
  \begin{itemize}
    \item Mientras la firma no invierte, posterga los flujos de caja que el
          proyecto ya estaría repartiendo. Ese costo de oportunidad es el
          \concepto{rendimiento por dividendos} \az{$\delta$}.
  \end{itemize}
  \cajaec{\az{\delta} = \az{\hat{\alpha}_v} - \rj{\alpha_v}
          \;\;\Longrightarrow\;\;
          \rj{\alpha_v} = \az{\hat{\alpha}_v - \delta} < \az{\hat{\alpha}_v}}
  \vspace{-2mm}
  \begin{ec}
    \text{Por lo tanto:}\quad \rj{\alpha_v} < \az{\hat{\alpha}_v}
  \end{ec}
  \vspace{-3mm}
  \begin{itemize}
    \item \az{$\hat{\alpha}_v$} es el retorno de equilibrio exigido a \rj{$V$};
          como parte llega como dividendo ($\az{\delta} > 0$), el valor latente
          se aprecia más lento.
    \item Si fuese $\az{\alpha_v} \geq \az{\hat{\alpha}_v}$, la opción
          \resultado{divergiría a infinito} y la firma jamás invertiría.
  \end{itemize}
\end{frame}
